\section{Giới thiệu}
\label{sec:introduction}

Các hệ thống hỏi--đáp pháp luật không chỉ cần tìm được văn bản có nội dung
liên quan mà còn phải xác định đúng điều kiện áp dụng, hệ quả pháp lý và căn
cứ điều luật của từng kết luận. Trong nhiều trường hợp, câu trả lời không
xuất hiện trực tiếp trong một điều khoản riêng lẻ mà phải được hình thành
bằng cách kết nối nhiều quy tắc. Một điều kiện ban đầu có thể làm phát sinh
một trạng thái pháp lý trung gian, sau đó trạng thái này tiếp tục kích hoạt
một hệ quả khác. Ngoài suy luận theo chiều thuận, người sử dụng còn có thể
yêu cầu truy ngược nguyên nhân, xác định sự kiện trung gian, kiểm tra một
quan hệ có phải là quan hệ trực tiếp hay không, hoặc đánh giá điều gì xảy ra
nếu một mắt xích trong chuỗi suy luận không còn tồn tại. Các benchmark pháp
luật như CaseHOLD và LegalBench đã cho thấy tính đa dạng và độ khó của các
tác vụ suy luận pháp lý, trong khi LegalBench-RAG nhấn mạnh thêm vai trò của
khâu truy hồi trong các hệ thống hỏi--đáp pháp luật
\cite{zheng2021casehold,guha2023legalbench,alami2024legalbenchrags}.

Retrieval-Augmented Generation (RAG) kết hợp mô hình sinh ngôn ngữ với nguồn
tri thức được truy hồi từ bên ngoài nhằm tăng tính xác thực và khả năng truy
vết của câu trả lời \cite{lewis2020rag}. Tuy nhiên, kiến trúc RAG truyền
thống thường biểu diễn corpus dưới dạng các đoạn văn tương đối độc lập và
xếp hạng chúng chủ yếu dựa trên độ tương đồng ngữ nghĩa với truy vấn. Cách
tiếp cận này có thể hoạt động tốt đối với câu hỏi một bước, nhưng gặp khó
khăn khi bằng chứng cần thiết được phân tán trên nhiều quy tắc hoặc khi thứ
tự liên kết giữa các bằng chứng quyết định ý nghĩa của câu trả lời. Các tác
vụ hỏi--đáp nhiều bước, tiêu biểu như HotpotQA, đã chỉ ra rằng việc truy hồi
đúng từng đoạn riêng lẻ chưa bảo đảm hệ thống khôi phục được toàn bộ đường
suy luận \cite{yang2018hotpotqa}.

Các phương pháp RAG dựa trên đồ thị giải quyết một phần hạn chế này bằng cách
biểu diễn thực thể, sự kiện hoặc đoạn tri thức dưới dạng node và quan hệ giữa
chúng dưới dạng edge. GraphRAG sử dụng cấu trúc đồ thị để hỗ trợ truy vấn và
tổng hợp thông tin ở phạm vi rộng, trong khi HippoRAG kết hợp truy hồi dựa
trên đồ thị với cơ chế liên kết nhiều bước
\cite{edge2024graphrag,gutierrez2024hipporag}. So với truy hồi các đoạn văn
độc lập, biểu diễn đồ thị giúp hệ thống giữ lại cấu trúc liên kết giữa các
mảnh bằng chứng. Tuy vậy, một đường đi bất kỳ trên đồ thị không nhất thiết
tương ứng với một chuỗi suy luận hợp lệ. Đối với miền pháp luật, hệ thống
còn phải phân biệt giữa quan hệ điều kiện--hệ quả được quy tắc hỗ trợ, quan
hệ trực tiếp, quan hệ thông qua sự kiện trung gian và đường thay thế dẫn đến
cùng một outcome.

Mô hình nhân quả cấu trúc (Structural Causal Model, SCM) cung cấp một khuôn
khổ hình thức để mô tả cơ chế sinh giá trị của các biến và phân tích kết quả
của can thiệp \cite{pearl2009causality}. Gần đây, một số nghiên cứu đã khảo
sát việc kết hợp causal reasoning hoặc counterfactual reasoning với RAG
\cite{wang2025causalrag,khadilkar2025causalcounterfactualrag,qin2026cfrag}.
Các hướng nghiên cứu này cho thấy tiềm năng của việc bổ sung cấu trúc nhân
quả vào quá trình truy hồi và sinh câu trả lời. Tuy nhiên, việc áp dụng các
khái niệm nhân quả vào pháp luật cần được giới hạn cẩn trọng. Quan hệ được
sử dụng trong nghiên cứu này là quan hệ điều kiện--hệ quả được thao tác hóa
từ quy tắc pháp luật, không phải quan hệ nhân quả được khám phá hoặc ước
lượng từ dữ liệu quan sát. Vì vậy, mô hình được xây dựng nên được hiểu là
một mô hình suy luận quy phạm có cấu trúc, phục vụ kiểm tra tính nhất quán
của chuỗi quy tắc và phân tích can thiệp trong phạm vi đồ thị đã cho.

Trên cơ sở đó, nghiên cứu này trình bày một pipeline CausalRAG cho hỏi--đáp
pháp luật hình sự Việt Nam. Pipeline bắt đầu bằng việc chuẩn hóa quy tắc
pháp luật thành các thành phần điều kiện, hệ quả, chủ thể và căn cứ điều
luật. Các điều kiện và hệ quả được ánh xạ thành sự kiện trong một đồ thị
nhân quả dị thể, trong đó các quy tắc đóng vai trò là cơ chế nối sự kiện
điều kiện với sự kiện hệ quả. Song song với đồ thị, hệ thống xây dựng bộ nhớ
vector cho rule và event để truy hồi các seed event liên quan đến truy vấn.
Từ các seed event, retriever mở rộng và xếp hạng các causal path có độ dài
tối đa hai hop, sau đó chuyển các đường dẫn ứng viên sang bước xác minh.

Khác với cách ra quyết định chỉ dựa trên số lượng path được hỗ trợ, thành
phần xác minh của chúng tôi phân tích nội dung cụ thể của truy vấn để xác
định loại claim đang được hỏi. Hệ thống phân biệt câu hỏi nhân quả chuẩn,
claim về cạnh trực tiếp, claim loại bỏ mediator, claim về tính cần thiết của
mediator và claim phản thực tế chưa xác định rõ. Việc phân tích này chỉ sử
dụng văn bản truy vấn và cấu trúc được truy hồi, không sử dụng
\texttt{question\_type}, câu trả lời chuẩn hoặc các nhãn đánh giá của
benchmark.

Để xác minh đường suy luận, chúng tôi xây dựng một mô hình nhân quả cấu trúc
pháp lý (Legal Structural Causal Model, LegalSCM). LegalSCM là một mô hình
suy luận tất định với miền trạng thái ba giá trị
$\{\mathrm{TRUE},\mathrm{FALSE},\mathrm{UNKNOWN}\}$. Đối với mỗi causal
path, mô hình kiểm tra liệu các cơ chế pháp lý tương ứng có thể kích hoạt
toàn bộ factual path hay không. Khi truy vấn yêu cầu phân tích phản thực tế,
LegalSCM thực hiện hard intervention $do(M=\mathrm{FALSE})$ trên mediator
$M$, cắt các cơ chế đầu vào của mediator và suy luận lại outcome. Nghiên cứu
đồng thời sử dụng một baseline đơn giản hơn, trong đó mediator được xóa khỏi
đồ thị và hệ thống kiểm tra liệu còn tồn tại đường thay thế đến outcome hay
không. Baseline này được gọi là \emph{path ablation}; nó phản ánh
reachability sau khi xóa node và không được diễn giải như một hard
intervention đầy đủ.

Kết quả xác minh được chuyển tới bước sinh câu trả lời theo intent. Thay vì
mặc định mọi câu hỏi đều yêu cầu hệ quả cuối cùng, hệ thống hỗ trợ nhiều dạng
đầu ra, bao gồm hệ quả cuối, nguyên nhân khởi đầu, sự kiện cầu nối, giải
thích từng bước, câu trả lời có/không, cấu trúc phân nhánh, cấu trúc hội tụ
và xác minh claim. Câu trả lời extractive được hình thành từ primary causal
path, ưu tiên nội dung hệ quả đầy đủ của quy tắc cuối và chỉ trích dẫn những
evidence thực sự được sử dụng.

Thực nghiệm được thực hiện trên corpus gồm 2.884 quy tắc và 1.814 sự kiện
được chuẩn hóa từ Bộ luật Hình sự Việt Nam. Đồ thị thu được chứa 5.134 node,
14.420 edge và 112 chuỗi nhân quả hai bước. Hệ thống được đánh giá trên một
benchmark tổng hợp gồm 250 câu hỏi, bao phủ suy luận thuận, suy luận ngược,
tìm sự kiện trung gian, giải thích chuỗi, câu hỏi có/không, phân nhánh và hội
tụ. Trên toàn bộ benchmark, hệ thống đạt Rule Recall@5 là $69{,}93\%$, Event
Recall@5 là $74{,}66\%$, độ chính xác xác minh là $92{,}00\%$, Answer Token
F1 là $68{,}69\%$ và Citation F1 là $62{,}86\%$. Khoảng cách giữa Top-1
Exact Path ($35{,}20\%$) và Oracle Exact Path ($71{,}20\%$) cho thấy nhiều
đường dẫn đúng đã xuất hiện trong tập ứng viên nhưng chưa được xếp ở vị trí
đầu tiên.

Paired ablation giữa LegalSCM và path ablation không ghi nhận khác biệt về
các quality metric trên benchmark hiện tại. Hai phương pháp đạt cùng độ
chính xác xác minh, Answer Token F1 và Citation F1. Tuy nhiên, LegalSCM cần
trung bình $3{,}890$ giây cho mỗi mẫu, so với $1{,}205$ giây của path
ablation, tương ứng chi phí thời gian cao hơn khoảng $3{,}23$ lần. Kết quả
này không cung cấp bằng chứng rằng structural verification cải thiện độ
chính xác trên benchmark hiện tại; thay vào đó, lợi ích quan sát được của
LegalSCM nằm ở semantics can thiệp rõ ràng hơn và khả năng tạo trace về
factual world, intervention assignment và counterfactual outcome.

Các đóng góp chính của nghiên cứu gồm:

\begin{itemize}
    \item Đề xuất một pipeline end-to-end kết hợp bộ nhớ vector, đồ thị nhân
    quả pháp luật, truy hồi nhiều bước, xác minh theo truy vấn và sinh câu
    trả lời có căn cứ điều luật.

    \item Xây dựng LegalSCM ba giá trị để kiểm tra factual causal path và
    thực hiện hard intervention trên sự kiện trung gian trong phạm vi các
    quy tắc pháp luật đã chuẩn hóa.

    \item Xây dựng một bộ đánh giá tổng hợp gồm 250 câu hỏi, bao phủ nhiều
    dạng suy luận pháp luật theo chiều thuận, chiều ngược, qua mediator,
    phân nhánh và hội tụ.

    \item Thực hiện paired ablation giữa structural SCM và node-deletion
    path ablation, qua đó xác định một null quality result trên benchmark
    hiện tại và định lượng chi phí thời gian của structural verification.

    \item Phân tích các giới hạn còn tồn tại, đặc biệt là lỗi xếp hạng causal
    path, hiệu năng thấp trên reverse reasoning và sự thiếu hụt các mẫu
    counterfactual có ground truth can thiệp.
\end{itemize}

Phần còn lại của bài báo được tổ chức như sau. Phần~\ref{sec:related-work}
trình bày các công trình liên quan về RAG, graph-based RAG, causal reasoning
và hỏi--đáp pháp luật. Phần~\ref{sec:problem} phát biểu bài toán và xác định
phạm vi tuyên bố. Phần~\ref{sec:method} mô tả pipeline và LegalSCM.
Phần~\ref{sec:data} trình bày corpus và benchmark. Thiết lập và kết quả thực
nghiệm lần lượt được trình bày trong Phần~\ref{sec:experimental-setup} và
Phần~\ref{sec:results}. Phần~\ref{sec:discussion} phân tích lỗi và thảo
luận kết quả. Cuối cùng, các hạn chế và kết luận được trình bày trong
Phần~\ref{sec:limitations} và Phần~\ref{sec:conclusion}.
